% Conclusion

\section{Conclusion}
\label{sec:conclusion}

We have presented a structure-preserving dynamical low-rank method for the
2D incompressible Navier--Stokes equations in stream-function form, and a
validation ladder that targets the regime where the interesting behavior
lives: forced, high-Reynolds-number turbulent dynamics at
$\mathrm{Re} \in \{100, 1000, 5000\}$. The scheme combines (i) an exact,
rank-preserving viscous step along the separable exponential flow of
Proposition~\ref{prop:viscous}, (ii) a second-order projected midpoint step
for the nonlinear part (eq.~\eqref{eq:step}), and (iii) a rank criterion
whose saturation at the grid's alias-free ceiling we measure and report
(Section~\ref{sec:rank-adaptation}). Divergence-freeness is exact by
representation, at every rank and every step (I1,
Table~\ref{tab:div}); the energy balance of eq.~\eqref{eq:energy} is
accounted for explicitly under forcing (I2).

The validation of Section~\ref{sec:results} is structured so that each
claim is checkable against the runs: the invariants I1--I4, the accuracy
against the full-grid spectral reference (Section~\ref{sec:res-error}), the
rank economy against the static POD baseline (Section~\ref{sec:res-pod}),
and the honest cost picture, including the regimes in which SP-DLRA is
slower than the full-grid reference (Section~\ref{sec:res-cost}).
% [PENDING-CODER: the headline numbers of the conclusion -- relative L2 error
% over the statistical window per Re, and the per-step and total cost ratios
% vs the full-grid reference and the POD baseline -- go here once the runs
% land; every number traceable to state/coder/results/ (CHECKLIST 1.1).]

Several directions follow directly from the open items identified in
Section~\ref{sec:limitations}. First, a three-dimensional formulation: the
exact viscous step carries over because $\Delta$ remains separable on the
periodic box, but the incompressibility constraint must then be carried by a
vector potential and gauge, and the nonlinear step must be redesigned around
the vortex-stretching term (Section~\ref{sec:disc-3d}). Second, a theory of
rank growth in forced turbulence. The measurement here is
that the retained rank does \emph{not} grow with Reynolds number: it reaches the budget within fifteen steps at every $\mathrm{Re}$ we
ran, and stays there, while the share of the energy in the zonal mode falls
from $20\%$ to $18\%$ across the same runs. Why a criterion driven by a
tolerance should report the resolved band rather than the flow, and what would
make it report the flow instead, is the question we think this leaves open.
% [PENDING-THEORETICAL-RESEARCH: if theoretical-research addresses why a
% tolerance-driven rank criterion reports the resolved band rather than the
% flow, cite it in Section 07 (sec:disc-rank) when finalizing this paragraph.]
Third, the forcing-aware invariant of the reduced model (Section~\ref{sec:limitations})
and a statistical-convergence study of the retained windows (V2). And fourth,
replacing the projected midpoint step with a robust (BUG-family) integrator,
as in the low-rank integrators of Ceruti et al.~\cite{ceruti2024}, which
offer better stability properties for stiff low-rank dynamics; we regard this
as a drop-in extension of the nonlinear step (eq.~\ref{eq:step}) and have
planned it as the next engine iteration.
